\chapter{Introduction}
Diffusion models, inspired by non-equilibrium statistical physics \cite{sohldickstein2015deepunsupervisedlearningusing}, are a popular paradigm of generative machine learning models that are prominently used for many different generation tasks such as image, video and audio synthesis as well as for scientific research tasks like 3D molecular generation and protein folding. \cite{ho2020denoisingdiffusionprobabilisticmodels} \cite{rombach2022highresolutionimagesynthesislatent} \cite{ho2022videodiffusionmodels} \cite{kong2021diffwaveversatilediffusionmodel} \cite{xu2022geodiffgeometricdiffusionmodel} More recently they have also been adopted heavily in world simulations and physical modeling, which is used in areas such as robotics \cite{yang2023unisimneuralclosedloopsensor}.

Training diffusion models means having them learn a distribution that matches the training data distribution as closely as possible. If the marginal distribution of the training data is somehow wrong, biased or doesn't reflect real world scenarios accurately then the generated samples of the trained model will reflect this in their output and give samples exhibiting this undesirable distribution. Even if the data distribution is fine there are still many situations where it is desirable to guide the generation process of the model toward some desired distribution, such as generating 50\% male and 50\% female images, or generating samples of particular colors.

Many different methods exist that allow for controlling the generation process. Incorporating the conditioning information into the training process \cite{dhariwal2021diffusionmodelsbeatgans} via classifier-guided diffusion or other similar conditional training approaches is one method but it requires training a task specific model along side large amounts of data, pairing the training data along the conditioning information. More flexible methods exist that allow for conditional sampling by directly working on the unconditional model. Twisted Diffusion Sampling \cite{wu2024practicalasymptoticallyexactconditional} and Feynman-Kac steering \cite{singhal2025a} are recent examples of approaches that do not require task-specific training and allow for flexible control of the sampling process. Normally, TDS is done when the generation is conditioned on a fixed, single observation $y$, where $y$ is defined through the likelyhood $p(y|x)$ \cite{wu2024practicalasymptoticallyexactconditional}, but combining it with a specific method for updating probability distributions, namely Jeffrey's rule, allows for steering toward a distribution instead. Feynman-Kac steering uses an arbitrary reward function to control the generation which naturally allows for steering toward a desired distribution.

Jeffrey's rule of conditioning is a generic way to update probability distributions in order to satisfy some given constraint. More specifically it says that the standard joint distribution
\begin{align*}
    p(z_I, z_J) = p(z_I | z_J)p(z_J)
\end{align*}
Can simply be rewritten as:
\begin{align*}
    p^*(z_I, z_J) = p(z_I | z_J)p^*(z_J)
\end{align*}
Where $p^*(z_J)$ is an adjusted marginal distribution of a subset of the features which had an undesirable marginal to begin with. This construction then leads to the formulation:
\begin{align*}
    p^*(z_I, z_J) = \frac{p^*(z_J)}{p(z_J)}p(z_I,z_J)
\end{align*}
Which says that the new joint distribution is an importance-weighted version of the original $p$. 
Combining Jeffrey's conditioning with the above mentioned methods for steering the sampling, it is possible to condition the sampling process on an arbitrary distribution. (CHECK? IS IT TRUE? FOR ANY ARBTRIARY DISTRIBUTION?)

This thesis explores using Jeffrey's rule CITATION needed(Jeffrey PDF from Jes? OR perhaps dig deeper into the refernces listed on there.) in combination with the algorithms for conditioning the sampling process of a diffusion model in order to steer the generation towards a desired target distribution. These methods will be evaluated and compared in terms of their sampling efficiency, distributional fidelity, and ability to maintain sample diversity while adhering to external constraints. (DOUBLE CHECK LAST SENTENCE, WHAT TO COMPARE?)
